% Setting Up, Chapter 3 -- Using the Keyboard (annex Sec 1: the SU part's
% orientation chapters; Sec 5 Conventions: key names in small capitals,
% hold-then-press and press-then-press; Sec 10 keyboard figure).
% Sources (Law 1): os/i123/app/face_a.c key() and command_key() (the keys
% Initech 123 acts on: the arrows, HOME, F2, F5, F9, ESC, BACKSPACE, ENTER,
% Ctrl-Break / Ctrl-C, / and <; a key it ignores does nothing), the shipped
% key map (harness/pubs/i123_face.c keymap: the probe finds exactly these
% keys acting; every \KeyRow below is graded by test-pubs-i123), the menu
% bar's File > Quit with its keyboard equivalent Ctrl-Q (face_a.c
% file_items), the templates (pubs/manuals/initech123-templates).
% GATES: test-flair-i123 test-i123-range-menu test-pubs-i123
\Gates{test-flair-i123, test-i123-range-menu, test-pubs-i123}
\chapter{Using the Keyboard}

You give Initech 123 almost all its instructions from the keyboard: you
type entries, move the cell pointer, and select commands from menus by
typing letters.  This chapter shows you which keys Initech 123 uses and how
the books in this package name them.

\section{The Parts of the Keyboard}
\ix{keyboard}
Your keyboard has four groups of keys that Initech 123 uses.
\begin{itemize}
\item The \textbf{typewriter keys} in the center: letters, digits and
  punctuation.  You type labels, numbers and formulas with them, and the
  first letters of menu commands.
\item The \textbf{pointer-movement keys}: the four arrow keys and
  \kname{HOME}.  They move the cell pointer around the worksheet, and the
  menu pointer along a menu.
\item The \textbf{function keys}, labelled F1, F2 and so on, across the
  top of the keyboard (101-key keyboard) or in two columns at its left
  (84-key keyboard).  Initech 123 gives three of them a function.
\item The \textbf{special keys}: \kname{ENTER}, \kname{ESC},
  \kname{BACKSPACE}, and the \key{Ctrl} key used with \key{Break}.
\end{itemize}

\section{Key Names in This Book}
\ix{key names}
This book prints the name of a key in small capitals: \kname{ENTER},
\kname{ESC}, \kname{HOME}.  A function key is named by what it does in
Initech 123, followed by the key itself in parentheses: \fkey{Goto}{F5}.
When you are to hold down one key while you press another, the two names
are joined by a hyphen: \kname{CTRL-BREAK} means hold down \key{Ctrl} and
press \key{Break}.  What you type is printed in a different typeface:
\type{@sum(b5..e5)}.

\section{The Keys Initech 123 Uses}
\ix{keys}
Besides the typewriter keys, these are all the keys Initech 123 acts on.
Any other key, such as \key{PgDn} or \key{F1}, does nothing in this
version.

\begin{KeyTable}{Key}{Name}
\KeyRow{UP}{Up arrow}{Moves the cell pointer up one row; in a menu, moves
  the menu pointer to the previous command.}
\KeyRow{DOWN}{Down arrow}{Moves the cell pointer down one row; in a menu,
  moves the menu pointer to the next command.}
\KeyRow{LEFT}{Left arrow}{Moves the cell pointer one column to the left;
  in a menu, to the previous command.}
\KeyRow{RIGHT}{Right arrow}{Moves the cell pointer one column to the
  right; in a menu, to the next command.}
\KeyRow{HOME}{Home}{Moves the cell pointer to A1; in a menu, to the first
  command.}
\KeyRow{F2}{EDIT}{Places the contents of the current cell on the control
  panel for changing.}
\KeyRow{F5}{GOTO}{Moves the cell pointer to the cell you name.}
\KeyRow{F9}{CALC}{Recalculates the formulas of the worksheet.}
\KeyRow{SLASH}{/}{Displays the main menu.}
\KeyRow{LESS-THAN}{<}{Displays the main menu, as \kname{SLASH} does.}
\KeyRow{ENTER}{Enter}{Completes an entry, selects the highlighted command,
  or accepts the answer to a prompt.}
\KeyRow{ESC}{Esc}{Cancels what you are doing and goes back one step.}
\KeyRow{BACKSPACE}{Backspace}{Erases the last character you typed.}
\KeyRow{CTRL-BREAK}{Ctrl-Break}{Leaves the menus, or a prompt, and returns
  to READY mode.}
\end{KeyTable}

The arrow keys move the cell pointer when the mode indicator shows
\ind{READY}; while you are typing an entry, an arrow key first completes
the entry, as \kname{ENTER} does, and then moves the cell pointer; and at
a prompt for a range they point to the cell you want.  \kname{HOME}
moves the cell pointer in \ind{READY} mode only, and the function keys
act in \ind{READY} mode only.

\begin{Note}
On a keyboard without a \key{Break} key, hold down \key{Ctrl} and press
\key{C}: it does the same as \kname{CTRL-BREAK}.
\end{Note}

\section{The Keyboard Templates}
\ix{keyboard templates}
Two keyboard templates come with Initech 123: a strip for the row of
function keys of the 101-key keyboard, and an L-shaped template for the two
columns of function keys of the 84-key keyboard.  Each template prints the
Initech 123 name of a function key above it.  A key with no name has no
function in Initech 123.  Place the template that fits your keyboard around
its function keys.

\section{The Keyboard and the FLAIR Menu Bar}
\ix{menu bar}
While Initech 123 is running, the menu bar at the top of the screen shows
its File menu, which has one command, \textbf{Quit}.  You can choose it
with the mouse, or press its keyboard equivalent \key{Ctrl-Q}.  It ends
Initech 123 in the same way as the worksheet command \slashcmd{/Quit Yes}
(Chapter~2).  All the other commands of Initech 123 are in its own menus,
which you display from the keyboard with \kname{SLASH}.

\section{Practicing}
\ix{practicing}
To get used to the keys before you begin the \doctitle{Tutorial}, start
Initech 123 and try these:
\begin{Keystrokes}
\Do{Press}{\kname{RIGHT} and \kname{DOWN} a few times}{to move the cell
  pointer; watch the address on the first line of the control panel}
\Do{Press}{\kname{HOME}}{to return to A1}
\Do{Press}{\kname{SLASH}}{to display the main menu; the mode indicator
  shows MENU}
\Do{Press}{\kname{RIGHT} several times}{to move the menu pointer along the
  menu; the third line describes each command}
\Do{Press}{\kname{ESC}}{to remove the menu}
\end{Keystrokes}
Nothing you do in this exercise changes the worksheet.
